

\section{Introduction}

Diffusion models~\cite{peebles2023dit} have achieved remarkable success in video generation; however, they face a severe computational bottleneck within the attention mechanism. Since video generative models often involve sequences exceeding $10\text{K}$ tokens, the attention cost scales quadratically with the token count. A promising approach to mitigate this is 4-bit quantization of FlashAttention, a method that is orthogonal to other strategies such as sparse attention\cite{zhang2025spargeattention,yang2025sparse,zhang2025sla,xi2025sparse,liu2021swin,chu2021twins,xiao2023efficient,xiao2024infllm,chen2024longlora,jiang2024minference,venkataramanan2024skipattention,yuan2024ditfastattn, zhang2025fast}. We develop a \textbf{highly compatible mixed-precision INT4 attention framework} for DiT-based video generation models with \textbf{3D-RoPE},
substantially improving quantization accuracy by selectively falling back to \textbf{FP16} in accuracy-sensitive attention blocks and denoising steps.

\paragraph{\textbf{Outliers and Rotation Transformation}}
The accuracy of low-bit quantized attention is constrained by the presence of \textit{outliers} in the tensors---elements whose magnitudes deviate significantly from the mean distribution. These extreme values compress the dynamic range available for the majority of the tensor, leading to severe precision loss. Rotation transformations are established techniques to mitigate outliers by redistributing tensor values, as discussed in detail in \cref{sec:relate}.


\paragraph{\textbf{Limitations of Existing Work.}}
Current 4-bit attention methods, when applied to DiT-based video generation models with 3D RoPE, face several limitations:
1) The Query ($\mathbf{Q}$) and Key ($\mathbf{K}$) matrices are usually not rotated, although rotation can effectively equalize outliers. This is because common rotation matrices are difficult to fuse with Rotary Position Embeddings (RoPE)~\cite{su2024roformer}, while online rotation introduces non-negligible overhead.
2) Existing methods employ symmetric quantization for the probability matrix $\mathbf{P} = \exp(\mathbf{S} - m)$. Since $\mathbf{P} \in (0, 1]$, this leaves half of the INT4 representation range unused.
Furthermore, LLM-oriented quantized attention methods with $\mathbf{Q}$/$\mathbf{K}$ rotations do not directly transfer to mixed-precision INT4 attention in 3D-RoPE video DiTs.




\paragraph{\textbf{Motivation and Key Findings.} }
In this work, we make the following observations: 
1) We investigate the impact of 3D RoPE on the outlier distribution of $\mathbf{Q}$ and $\mathbf{K}$ in 3D-RoPE video DiTs. We identify specific patterns in the outlier distribution that can be exploited to implement a highly efficient and effective rotation matrix. 
2) Fully utilizing the 4-bit numerical range for the quantization of $\mathbf{P}$ yields a significant improvement in attention accuracy. 

\paragraph{\textbf{Method.} }
We introduce \textbf{RotateAttention}, which comprises two core innovations complemented by a standard offline technique: 
1) \textbf{RoPE-aware Rotation:} A block-diagonal matrix composed of $2 \times 2$ orthogonal blocks that effectively mitigates outliers in $\mathbf{Q}$ and $\mathbf{K}$ with minimal overhead, designed to be compatible with or mergeable into RoPE. 
2) \textbf{Range-Optimized $\mathbf{P}$ Quantization:} A fixed-scale and zero-point calibration for matrix $\mathbf{P}$ that fully exploits the INT4 dynamic range at low computational cost. 
Additionally, we apply a \textbf{Hadamard Transform for $\mathbf{V}$}, a well-established technique~\cite{liu2025spinquant,sun2025flatquant} that mitigates outliers in $\mathbf{V}$ by fusing Hadamard matrices into the existing linear projections at zero inference cost.

Experimental results demonstrate that our method is both efficient and effective. More importantly, our strategy can be integrated into other quantized attention frameworks to further enhance their performance.


